\chapter{Formal Sudoku Rules}

\section{State}\label{sec:state}

\paragraph{Size}
Size of an $n \times n$ sudoku is a positive integer $n$.

\paragraph{State}
State $S := \{S_{i,j} | i, j \in [n]\}$.
Each element $S_{i,j} \subseteq [n]$.

\subsection{Task}

Task $T := \{T_{i,j} | i, j \in [n]\}$ is the initial state of an unsolved sudoku puzzle.
Each element $T_{i,j}$ represents what numbers can be in the cell $(i, j)$ of a sudoku puzzle.

\subsection{Solution}

Solution is a sudoku state, such that every element $S_i$ has exactly one element.
It represents solved sudoku puzzle.
Each element $S_{i, j}$ represents what number is in the cell $(i, j)$ of a solved sudoku puzzle.

Solution $S$ is a valid solution of a sudoku puzzle with task $T$, if $S_{i,j} \subseteq T_{i,j}$ for every $(i, j)$.

\section{Rules}

Section~\ref{sec:state} only requires valid solution to choose one number for each cell from task.
Sudoku rules further restrict what a valid solution can be.

\subsection*{Regions}

Let us have a sudoku puzzle with task $T$ and solution $S$.

Sudoku puzzle has a finite set of regions $R := {R_1, R_2, ..., R_r}$.
Each region $R_i$ is a set of n coordinates of cells of the puzzle.
\[
    \left|R_i\right| = n \land R_i \subseteq \{(i, j) | i, j \in [n]\}
\]

Each region $R_i$ adds an extra requirement on a valid solution of the puzzle:
Cells in a region must contain each number exactly once.

\subsection*{Variants}

Variant can add extra requirements for a valid solution, that can not be encoded by regions.

A variant only needs to contain one validation function $v(S)$,
which takes a puzzle's solution and returns true if the solution does not break any rule of the variant
and returns false it it does.
If validation function returns false, the solution can not be valid.

\subsection*{Edge variant}

Edge variant is a specific type of variant.
All variants implemented in the software are edge variants.

Edge variant $E := (L, E, E_L, L_N, V)$.

$L := [l]$
is a set of labels.

$E \subseteq \{(a, b) | a, b \in [n] \times [n] \land a \text{ is edge-adjacent to } b\}$
is a set of edge-adjacent cells with a label.

$E_L : E \to L$
is an edge-label function, that for an edge-adjacent pair of cells with a label returns the label.

$L_N : L \to 2^{[n] \times [n]}$
is a label-neighbourhood function, that for a label returns set of all (oriented) allowed pairs of numbers for that label.

$V$ is a validation function that returns true, if for each pair of edge-adjacent cells $(a, b) \in E$ with a label $E_L((a, b))$:
The pair of numbers $\{S_a, S_b\}$ in solution $S$ is allowed by the label-neighbourhood function for the label.

In other words: $V(S) := \bigwedge_{(a, b) \in E} (S_a, S_b) \in L_N(E_L((a, b)))$
